Distribution-independent statistical-query learning does not imply low dimension complexity, even for nondegenerate probabilistic variants. If every product distribution has a monochromatic rectangle of mass at least $1/\rho$, a common learner uses $O(\rho(\log|H|+\log(1/\epsilon)))$ queries at tolerance $\Omega(\epsilon/(\rho\log(1/\epsilon)))$. For random incidence flips on $2M$ points, this yields $O_\epsilon(\log M)$ queries despite dimension $\Omega(M^{1/3}/\log M)$, exact probabilistic dimension $\Omega(M^{1/3}/\log^2 M)$, and averaged probabilistic dimension $\Omega_\eta(M^{1/3}/\log M)$ for fixed $\eta<1/2$. Proper deterministic implementations run in time polynomial in the supplied table. Under subexponential pseudorandom-function security, a key-known polynomial-time learner uses $O_\epsilon(n^2)$ queries while almost every key gives dimension at least $n^c/(30\log_2 n)$, for any prescribed fixed $c$ and $n=\log_2|X|$. These classes also admit engineered polynomial-size differentiable-learning simulations. Bounded product-average communication is sufficient for SQ learning, but cube halfspaces show it is unnecessary. The prescribed Gaussian-initialized, bias-free online-SGD question remains open.
